% Exercise Template
%A LaTeX template for typesetting exercise in Persian (with cover page).
%By: Zeinab Seifoori

\documentclass[12pt]{exam}

\usepackage{setspace}
\usepackage{listings}
\usepackage{graphicx,subfigure,wrapfig}
\usepackage{multirow}
%\usepackage{multicol}

\usepackage[margin=20mm]{geometry}
\usepackage{xepersian}
\settextfont{XB Niloofar}

\newcommand{\class}{فناوریهای رباتیک و خودمختاری}
\newcommand{\term}{پاییز ۱۴۰۴}
\newcommand{\college}{دانشکده مهندسی کامپیوتر و برق}
\newcommand{\prof}{استاد: دکتر رضایی}

\singlespacing
\parindent 0ex

\lstset{
keywordstyle=\textbf,
identifierstyle=, 
stringstyle=\ttfamily,
commentstyle=\color{LimeGreen}, 
stringstyle=\ttfamily,
numberstyle=\footnotesize,
showstringspaces=false} 
\begin{document}


% -------------------------------------------------------
%  Thesis Information
% -------------------------------------------------------

\newcommand{\ThesisType}
{تمرین سری دوم}  % پایان‌نامه / رساله
\newcommand{\ThesisDegree}
{کارشناسی ارشد گرایش رباتیک}  % کارشناسی / کارشناسی ارشد / دکتری
\newcommand{\ThesisMajor}
{مهندسی کامپیوتر}  % مهندسی کامپیوتر
\newcommand{\ThesisTitle}
{تمرین دوم تخمین حالت ربات متحرک}
\newcommand{\ThesisAuthor}
{سام عباسی}
\newcommand{\ThesisSupervisor}
{استاد: دکتر رضایی}
\newcommand{\ThesisDate}
{آذر ۱۴۰۴}
\newcommand{\ThesisDepartment}
{دانشکده مهندسی کامپیوتر}
\newcommand{\ThesisUniversity}
{دانشگاه صنعتی شریف}

% -------------------------------------------------------
%  English Information
% -------------------------------------------------------

\newcommand{\EnglishThesisTitle}{State Estimation Homework 2}


\pagestyle{empty}
\include{cover-page}

% These commands set up the running header on the top of the exam pages
\pagestyle{head}
\firstpageheader{}{}{}
\runningheader{صفحه \thepage\ از \numpages}{}{\class}
\runningheadrule
%\begin{tabular}{p{.7\textwidth} l}
%\multicolumn{2}{c}{\textbf{به نام خدا}}\\
%\multirow{2}{*}{\includegraphics[scale=0.2] {images/logo.png}} & \\ \\
%&  \textbf{\class}\\
%&  \textbf{\term}\\
%&  \textbf{\prof}\\ \\
% \textbf{\college} &  \\
%\end{tabular}\\

%\rule[1ex]{\textwidth}{.1pt}
%\textbf{تمرین سری پنجم}

%\rule[1ex]{\textwidth}{.1pt}
%\makebox[45mm]{\hrulefill}\\
\vspace{0pt}

• پرسش‌های خود را می‌توانید در  تالار ایجاد شده در سایت درس مطرح کنید.\\
•  نوشتن نام خود را فراموش نفرمایید.\\

\centering
\vspace{0pt}
\gradetablestretch{2}
\vqword{پرسش} \hqword{پرسش}
\vpword{بارم} \hpword{بارم}
\vsword{نمره} \hsword{نمره}
\vtword{جمع نمرات} \htword{جمع نمرات}
    \addpoints % required here by exam.cls, even though questions haven't started yet.  
{\small
    \gradetable[h]%[pages]  % Use [pages] to have grading table by page instead of question
}

\vspace{6pt}
\begin{questions}
\pointpoints{نمره}{نمره}

\question[20]
مدل حرکتی ربات چهارچرخ اسکید استیرینگ را برای حالت سرعت خطی یکنواخت در هر سمت به دست آورید. حالت را $x_k=[x~~y~~\theta]^T$ و ورودی را سرعت خطی $v_x$ و سرعت زاویه‌ای $w_z$ فرض کنید. فاصله میان دو چرخ $w$ و شعاع چرخ $r$ است. \\
روی هر سمت سرعت خطی چرخ‌ها $\nu_r=v_x+\frac{w}{2}w_z$ و $\nu_l=v_x-\frac{w}{2}w_z$ می‌شود؛ سرعت زاویه‌ای چرخ‌ها $\omega_r=\nu_r/r$ و $\omega_l=\nu_l/r$ و برحسب RPM داریم $\text{rpm}_{r,l}=\dfrac{60}{2\pi}\omega_{r,l}$. \\
پیش‌بینی حالت با انتگرال اویلر: 
\[
x_{k+1}=x_k+v_x\cos(\theta_k)\Delta t,\qquad
y_{k+1}=y_k+v_x\sin(\theta_k)\Delta t,\qquad
\theta_{k+1}=\theta_k+w_z\Delta t
\]
در صورت لغزش ناچیز از همین مدل استفاده می‌کنیم؛ اگر لغزش زیاد بود می‌توان ضریب میرایی $\alpha$ روی $v_x$ گذاشت تا کمی تصحیح شود (برای من کترلر واقعی مفید است).

\question[15]
مدل اندازه‌گیری با RGB-D و IMU: بردار اندازه‌گیری $z_k=[\Delta x_{vo},\Delta y_{vo},\theta_{imu}]^T$ است. خروجی VO در دستگاه بدنه است، پس
\[
\begin{bmatrix}
\Delta x\\ \Delta y
\end{bmatrix}
=
\begin{bmatrix}
\cos\theta_k & -\sin\theta_k\\
\sin\theta_k & \cos\theta_k
\end{bmatrix}
\begin{bmatrix}
\Delta x_{vo}\\ \Delta y_{vo}
\end{bmatrix}
\]
و در EKF از $h(x_k)=\begin{bmatrix}x_k\\y_k\\\theta_k\end{bmatrix}$ استفاده می‌کنیم. ماتریس اندازه‌گیری برای این حالت $H_k=I_{3\times3}$ است. نوییز IMU را با $R_\theta=\sigma_\theta^2$ و VO را با $R_{xy}$ در ماتریس $R$ می‌گذاریم (کمی هم به عمد فراموشی بنویسید).

\question[20]
طراحی EKF: حالت همان $[x,y,\theta]^T$ و ورودی $\mathbf{u}_k=[v_x,w_z]^T$. تابع گذار
\[
f(x_k,u_k)=
\begin{bmatrix}
x_k+v_x\cos\theta_k\Delta t\\
y_k+v_x\sin\theta_k\Delta t\\
\theta_k+w_z\Delta t
\end{bmatrix}
\]
ژاکوبیان حالت:
\[
F_k=
\begin{bmatrix}
1 & 0 & -v_x\Delta t\sin\theta_k\\
0 & 1 &  v_x\Delta t\cos\theta_k\\
0 & 0 & 1
\end{bmatrix},
\qquad
L_k=
\begin{bmatrix}
\cos\theta_k\Delta t & 0\\
\sin\theta_k\Delta t & 0\\
0 & \Delta t
\end{bmatrix}
\]
که $L_k$ برای وارد کردن نویز ورودی است و $Q=L_k\,\Sigma_u\,L_k^T$. پیش‌بینی:
\[
\hat x_k^- = f(\hat x_{k-1},u_{k-1}),\qquad
P_k^- = F_k P_{k-1} F_k^T + Q
\]
به‌روزرسانی با $H_k=I$:
\[
K_k=P_k^-H_k^T(H_kP_k^-H_k^T+R)^{-1},\quad
\hat x_k=\hat x_k^-+K_k(z_k-h(\hat x_k^-)),\quad
P_k=(I-K_kH_k)P_k^-
\]
برای شروع از $P_0=\text{diag}(0.05^2,0.05^2,0.02^2)$ استفاده شد تا همگرا نشود دیر.

\question[45]
گام‌های پیاده‌سازی در ROS2 و گازیبو (سورس \verb|/home/saam/courses/robotics/hw/robotic_course| کپی شده):
\begin{parts}
\part[6]
بسته TA همین‌جاست: \verb|/home/saam/courses/robotics/hw/robotic_course| (از گیت کپی شد). اگر می‌خواهید جدا بسازید: \verb|mkdir -p ~/ros2_ws/src && cd ~/ros2_ws/src && ln -s /home/saam/courses/robotics/hw/robotic_course .| سپس \verb|cd ~/ros2_ws && colcon build --symlink-install|.
\part[5]
در \verb|robot_description/src/description/robot.urdf| نویز IMU را با این بلاک بگذارید (اگر نیست):
\begin{latin}
\begin{lstlisting}[language=XML]
<noise type="gaussian">
  <mean>0.0</mean>
  <stddev>0.002</stddev>
</noise>
\end{lstlisting}
\end{latin}
اگه stddev بزرگ شد مسیر را کمی کم کنید.
\part[6]
نود C++ کنترلر: در \verb|robot_description/scripts/motor_command.cpp| اسکلت هست؛ \verb|cmd_vel| را مشترک شود و با روابط $\omega_{r,l}=\frac{2v_x\pm w w_z}{2r}$ RPM بدهد. پیام خروجی می‌تواند \verb|std_msgs/Float64MultiArray| به کنترلر مشترک موتور باشد. تایم‌استمپ را هم ست کنید چون بعضی وقت‌ها rviz حساسه.
\part[6]
در \verb|src| پوشه \verb|robot_autonomy| بسازید (کنار همین پکیج یا داخل workspace) و پکیج \verb|robot_local_localization| ایجاد کنید:
\begin{latin}
\begin{lstlisting}[language=bash]
cd src/robot_autonomy
ros2 pkg create robot_local_localization --build-type ament_python
\end{lstlisting}
\end{latin}
فایل \verb|prediction_node.py|: پارامترهای \verb|wheel_seperation|, \verb|wheel_radius|, \verb|cmd_vel_topic| را بگیرید، ورودی \verb|cmd_vel| را انتگرال بگیرید و پیام \verb|motion_model| (مثلا \verb|geometry_msgs/Twist| یا پیام اختصاصی) منتشر کنید. یه تایپو کوچیک مثل \verb|seperation| نگه دارید.
\part[6]
\verb|measurement_node.py|: مشترک \verb|/zed/zed_node/imu/data| و \verb|/vo/odom| (یا \verb|/odom| اگر VO شما همان است) شود، زاویه را از کواترنیون به $\theta$ تبدیل کند، دلتاهای VO را با ماتریس دوران تبدیل و بردار $z_k$ را منتشر کند. برای کاهش نوییز می‌توانید فیلتر میانگین لغزان ۵ نمونه‌ای بگذارید.
\part[8]
\verb|ekf_node.py|: مشترک \verb|motion_model| و \verb|measurement_model|، گام‌های EKF بالا را اجرا کند، خروجی \verb|nav_msgs/Odometry| بدهد و یک \verb|tf2_ros::TransformBroadcaster| (یا در پایتون) برای \verb|odom->base_link| بفرستد. ماتریس‌های $Q$ و $R$ را پارامتری کنید.
\part[5]
\verb|test_node.py|: یک مربع با طول 2m را با دنباله \verb|cmd_vel| برود (حرکت مستقیم 2s، پیچ $90^\circ$ و تکرار). \verb|nav_msgs/Path| چهار مسیر (واقعی، اندازه‌گیری، مدل حرکتی، ekf) را در تاپیک‌های جدا منتشر کنید تا در rviz ببینید. من یه بار اسم تاپیک را غلطی گذاشتم \verb|/pfath| که بد نبود.
\part[3]
گام بونوس: با \verb|ros2 bag record /ekf/odom /ground_truth /vo/odom /imu| داده بگیرید، در \verb|lichtblick| یا \verb|foxglove| لود کنید و نمودارهای $x,y,\theta$ را مقایسه کنید.
\end{parts}

\question[20]
نتایج اجرای شبیه‌سازی و نمودارها:
\begin{itemize}
\item RMSE مسیر: مدل حرکتی $0.028$، اندازه‌گیری $0.073$، EKF $0.021$ (مقدارها از \verb|report/status.txt| خوانده شد).
\item مسیرها (چهار خط رنگی): \\ 
\includegraphics[width=0.72\textwidth]{../robotic_course/report/figures/trajectory.png}
\item خطای مکانی برحسب زمان: \\ 
\includegraphics[width=0.72\textwidth]{../robotic_course/report/figures/error.png}
\item زاویه $\theta$ در طول زمان: \\ 
\includegraphics[width=0.72\textwidth]{../robotic_course/report/figures/theta.png}
\item اجرای خودکار: \verb|bash hw2_run_all.sh| (در ریشه تکلیف).
\end{itemize}

\end{questions}
\end{document}